% !TEX root = ../../Main.tex
\section{Benchmarks and Performance Metrics}

The benchmark is to buy the pair at the start of the window and hold it to the end. This is the fair comparison for a directional strategy on a single instrument. Anyone can do it, and it costs one spread. A strategy that cannot beat it does not justify its extra complexity. The thesis proposal named a random-signal baseline and a moving-average crossover instead. Both were dropped, because a trader would not hold either of them as a real alternative.

The reported metrics are those defined in \Cref{Section:PerformanceMeasurement}: total and annualised return, the Sharpe, Sortino, Calmar and Sterling ratios, maximum drawdown and the number of trades. \Cref{Tables:PerformanceMetrics} states what each one measures. All ratios use a risk-free rate of zero. They can therefore be compared across pairs without any assumption about interest rates.

% !TEX root = ../Main.tex
\begin{table}[htb!]
\centering
\footnotesize
\caption{Summary of the performance metrics used in this work.}
\label{Tables:PerformanceMetrics}
\begin{tabularx}{\textwidth}{@{}lX@{}}
\toprule
\textbf{Metric} & \textbf{Description} \\
\midrule
Net Profit & Total profit or loss \\
\addlinespace
Profit Factor & Ratio of gross profit to gross loss \\
\addlinespace
Commission and Swap & Costs that reduce net profit \\
\addlinespace
Max Balance and Equity Drawdown & Largest losses observed \\
\addlinespace
Total Trades and Winning Trades & How much the strategy trades and how often it wins \\
\addlinespace
Max Consecutive Winning and Losing Trades & How consistent the results are \\
\addlinespace
Largest Winning and Losing Trades & How large single gains or losses can be \\
\addlinespace
Average Trade & Average result per trade, a measure of trading efficiency \\
\addlinespace
Sharpe Ratio and Sortino Ratio & Risk-adjusted return \\
\bottomrule
\end{tabularx}
\end{table}

For four of them, it is worth saying why all four are reported instead of choosing one. Return says how much was made. It says nothing about the risk taken to make it. The Sharpe ratio divides return by total volatility. This penalises a strategy for moving up as much as for moving down. The Sortino ratio corrects this by counting only downside movement. Calmar and Sterling divide by drawdown instead of volatility. Calmar uses the worst loss ever suffered and Sterling the average one. They therefore measure how far the account fell, and not how volatile it was, which is the question an investor asks. A strategy can look strong on one of these and weak on another. Where that happens in this work, the text points it out.

One point applies to every number in this chapter. Returns are net of the spread quoted at the moment of each trade and of the account commission in \Cref{Tables:CostModel}. They therefore cannot be compared with results computed on bar closes with a fixed spread or with no spread. They are also lower than the same policies would show under those conditions.